\documentclass{article}
\title{\textbf{Trabajo Práctico Especial - \rplang{}} \\ [1ex]
\large Instituto Tecnológico de Buenos Aires - Autómatas, Teoría de Lenguajes y Compiladores (72.39) \\ [1ex]
\large Grupo RP-Lang}
\date{21 de junio de 2025}
\author{
\textbf{Ignacio Searles}\\
isearles@itba.edu.ar\\
64.536
\and
\textbf{Augusto Barthelemy Solá}\\
abarthelemysola@itba.edu.ar\\
64.502
\and
\textbf{Santiago Bassi}\\
sabassi@itba.edu.ar\\
64.643
}

\usepackage{multicol}
\usepackage{graphicx, wrapfig}
\graphicspath{ {imágenes/} }

\usepackage{float}
\usepackage{amsmath}
\usepackage{amsfonts}

\usepackage{hyperref}

\usepackage{caption, threeparttable}
\usepackage{hyperref}
\usepackage{listings}

\usepackage[margin=1.3in]{geometry}
\newcommand{\rplang}{\texttt{RP-Lang}}

\renewcommand{\figurename}{Figura}
\renewcommand{\tablename}{Tabla}
\renewcommand{\abstractname}{Resumen}

\begin{document}
\maketitle

\begin {abstract}

El presente informe trata sobre el diseño e implementación de un lenguaje que permite la definición y evaluación de funciones recursivas primitivas. El compilador de \rplang{} traduce programas escritos en este lenguaje a código C, aprovechando la expresividad y eficiencia de dicho lenguaje de bajo nivel para ejecutar las funciones definidas. 

\end {abstract}

\tableofcontents

\section {Introducción}

Las funciones recursivas primitivas constituyen una clase fundamental de funciones total computables, ampliamente estudiadas en teoría de la computación. En este trabajo se presenta \rplang{}, un lenguaje de programación diseñado para expresar y evaluar exclusivamente este tipo de funciones.

El desarrollo del lenguaje incluye el diseño de su sintaxis y semántica, así como la implementación de un compilador que traduce programas \rplang{} a código en lenguaje C. Para ello, se utilizaron las herramientas Flex y Bison, facilitando la construcción de los analizadores léxico y sintáctico respectivamente.

El objetivo principal del proyecto es proporcionar una herramienta que permita experimentar con funciones recursivas primitivas de forma práctica.

\section {Modelo computacional}

\subsection {Dominio}

Desarrollar un lenguaje que permita definir funciones recursivas primitivas (RP). Las funciones RP son de la forma $f:\mathbb N^k\to\mathbb N$ (Zahnd, 2022). Las mismas cumplen que están definidas en base a los esquemas recursivos de tipo I y tipo II y composición a partir de las funciones iniciales:
\begin{itemize}
\item $suc:\mathbb N \to \mathbb N, suc(x) = x + 1$
\item $const_k:\mathbb N \to \mathbb N, const_k(x) = k$
\item $\pi^k_j:\mathbb N^k \to \mathbb N, \pi^k_j(x_1, \cdots, x_k) = x_j$
\end{itemize}

Las definiciones de funciones RP pueden ser composiciones o esquemas recursivos:
\begin{itemize}
\item Para la composición, si $h$ y $g$ son RP $\implies f(x) = g(h(x))$ es RP. 
\item Para el esquema recursivo de tipo I, si $g:\mathbb N^2\to\mathbb N$ es RP $\implies f:\mathbb N\to\mathbb N$ tal que $f(0) = k \in \mathbb N$ y $f(n + 1) = g(n, f(n))$ es RP.
\item Para el esquema recursivo de tipo II, $g:\mathbb N^{n+2}\to\mathbb N$ y $q:\mathbb N^{n}\to\mathbb N$ son RP $\implies f:\mathbb N^{n+1}\to\mathbb N$ tal que $f(x_1, \cdots, x_n, 0) = q(x_1, \cdots, x_n)$ y $f(x_1, \cdots, x_n, y + 1) = g(x_1, \cdots, x_n, y, f(x_1, \cdots, x_n, y))$ es RP.
\end{itemize}

El lenguaje permite evaluar cualquier función RP que se haya definido. Haciendo uso de las definiciones de las mismas para evaluarlas recursivamente y obtener un valor de salida. 

\subsection {Lenguaje}

Un programa de \rplang{} (extensión recomendada $\texttt{.rp}$), está compuesto por definiciones, evaluaciones e $\texttt{\#include}$. Los mismos puede aparecer intercalados en cualquier orden, mientras se cumpla que no se utilicen funciones no definidas en una definición. A continuación se listan construcciones del lenguaje:

\begin{itemize}
\item Se pueden definir funciones partiendo de las funciones básicas.
\item Se pueden definir funciones mediante composición de funciones.
\item Se pueden definir funciones mediante esquemas recursivos de tipo I y II.
\item Se pueden evaluar funciones utilizando constantes o entradas por teclado.
\item Se pueden importar funciones desde otros archivos.
\item Las funciones binarias se pueden evaluar mediante la syntaxis: $(\texttt{expr1}\ \texttt{func}\ \texttt{expr2})$ donde $f:\mathbb N^2\to\mathbb N$.
\item Se pueden incluir comentarios en el código.
\end{itemize}

Las funciones pueden definirse mediante expresiones compuestas o mediante esquemas recursivos. A continuación se muestran ejemplos representativos de cada caso.

\paragraph{Composición}

\begin{verbatim}
#include times.rp

def pow3(x):
=> pow3(x) = (x * (x * x))

//this will prompt the user for input
pow3(input)
\end{verbatim}

\paragraph{Recursión (esquema tipo I)}

\begin{verbatim}
#include times.rp

def !(n):
-> !(0) = 1
-> !(n + 1) = *(suc(n), !(n))

!(4)
\end{verbatim}

\paragraph{Recursión (esquema tipo II)}

\begin{verbatim}
#include times.rp

def ^(x, y):
-> ^(x, 0) = 1
-> ^(x, y + 1) = *(x, ^(x, y))

(2 ^ 5)
\end{verbatim}

\section {Implementación}
\subsection {Frontend}
\subsection {Backend}
\subsection {Dificultades encontradas}

\section {Futuras Extensiones}
\section {Conclusiones}
\section {Referencias}
\section {Bibliografía}

A continuación, se enumeran las fuentes consultadas para la elaboración de este trabajo.
\begin {itemize}
\item \href{https://github.com/mzahnd/apuntes-9335-logica-computacional}{Zahnd, M. (2022). GitHub - mzahnd/apuntes-9335-logica-computacional: Apuntes de la materia Lógica Computacional (93.35) del ITBA. GitHub.}
\end {itemize}

\end{document}
